\section{Reproducibility Statement}
\label{sec:reproducibility}
% 【中文】\section{可复现性声明}

We release a \emph{protocol white-paper} (harness protocol, evaluation criteria, desensitized task specifications) and an \emph{NDA-controlled review package} granting reviewers access to a sandbox with the commercial EDA tools. We do \textbf{not} release executable Lite-set assets or a mock-EDA mode (v1.4 open-source boundary narrowing). The falsifiability of our claims does not depend on proprietary licenses: an independent party can re-derive an equivalent benchmark on open-source EDA tools (Yosys/OpenROAD/Magic) and test the same 3 core + 7 sub-hypotheses using the released protocol. The M0 baseline (\MzeroPass{}) is reproducible from the frozen 696-run snapshot configuration (DeepSeek-V4-Pro, Cline v3.0.51, \texttt{thinking=low}, 2400s timeout). Full benchmark construction, per-domain results, and experimental details are in Appendices~\ref{sec:benchmark-details}--\ref{sec:experimental-details}.
% 【中文】
% 我们发布\emph{协议白皮书}（执行框架协议、评测判据、脱敏任务规格）与\emph{NDA 受控评审包}，授予
% 评审者访问含商用 EDA 工具的沙箱。我们\textbf{不}发布可执行的 Lite 集资产或 mock-EDA 模式
%（v1.4 开源边界收缩）。我们主张的可证伪性不依赖于专有许可：独立方可在开源 EDA 工具
%（Yosys/OpenROAD/Magic）上重新推导出等价基准，并用发布的协议检验同样的 3 条核心 + 7 条子假设。
% M0 基线（\MzeroPass{}）可从冻结的 696-run 快照配置复现（DeepSeek-V4-Pro、Cline v3.0.51、
% \texttt{thinking=low}、2400s 超时）。完整基准构建、各领域结果与实验细节见
% 附录~\ref{sec:benchmark-details}--\ref{sec:experimental-details}。
